\documentclass{article}
\title{時系列分析入門 — 偶然の時間構造を読む}
\author{TeX 図書館}

\begin{document}

\section{この教科書のために — 時間を持つデータ}

これまで統計の主役は「独立なサンプル」でした。しかし価格・気温・売上のデータは\textbf{昨日の値と今日の値がつながっている}——独立という前提が崩れる世界です。時系列分析は「\textbf{時間構造をもつ偶然}」を扱う統計の分野で、『確率過程入門』(理論)と実データの間の橋です。

\section{定常性 — いちばん大事な概念}

\textbf{定常}( stationary )な時系列とは、\textbf{平均も分散も自己相関構造も時間によらない}系列です。定常なら過去の統計が未来で使えます。非定常(トレンドを持つなど)なら、過去の平均は未来の予測として無意味です。

\begin{quote}
\textbf{【例】}\ 「日々の気温」は季節性があるため非定常。「気温の\textbf{前日比}」はほぼ定常。株価そのものは非定常(水準が伸びる)ですが、\textbf{対数リターン(前日比)}はほぼ定常とみなせる——\textbf{差分をとる}ことが非定常を定常に変える基本操作です。
\end{quote}

\begin{quote}
\textbf{【演習】}\ 「銀行口座残高」と「毎月の支出」のうち、どちらが定常に近いか、また理由を述べよ。
\end{quote}

\textbf{解答の指針:}\ 支出(前月比の変動は安定しがち)。残高は積み上がるため非定常(トレンド)。

\section{自己相関 — ラグ構造を測る}

\textbf{ラグ \(k\) の自己相関係数}( ACF ):
\[ \rho_k = \frac{\mathrm{Cov}(X_t, X_{t-k})}{V[X_t]} \]
「\(k\) 期前の自分と、どれだけ一緒に動くか」。定常系列では \(\rho_1, \rho_2, \dots\) を計算すると、系列の\textbf{記憶の長さ}が読めます:
\begin{itemize}
\item すぐ 0 に落ちる → 記憶が短い(ホワイトノイズ)
\item ゆっくり減衰する → 強い持続性(トレンド・長期記憶)
\item 周期をもつ → 季節性
\end{itemize}

\begin{quote}
\textbf{【注意:平滑化の自己相関】}\ 移動平均をかけた系列は\textbf{人工的に自己相関が高くなります}(隣接項が共通のデータを含むため)。「移動平均後に相関が見つかった」分析は偽の構造を拾いがちです。
\end{quote}

\begin{quote}
\textbf{【演習】}\ ホワイトノイズ(独立な系列)の ACF の理論的な値を、すべてのラグで述べよ。
\end{quote}

\textbf{解答の指針:}\ 理論値はラグ 0 のみ 1、他はすべて 0。標本では標準誤差約 \(1/\sqrt{n}\) の範囲で揺れる。

\section{ランダムウォーク検定 — 「予測できない」を確かめる}

価格系列に\textbf{ランダムウォーク仮説}(前日比が予測不可能)が成立するかを検定するのは、実務上極めて重要です。

\begin{enumerate}
\item \textbf{ACF 検定}: リターン系列の自己相関が有意に 0 から離れているか。ラグ 1〜10 でほぼ有意なら「何らかの予測可能性」の存在
\item \textbf{ラン(chi)検定}: 上昇/下落の連(ラン)の長さの分布が、独立なコイントスと比べて長いか短いか
\item \textbf{分散比検定}: \(n\) 日リターンの分散が 1 日リターンの分散の \(n\) 倍(ランダムウォークの理論値)から逸脱していないか
\end{enumerate}

\begin{lstlisting}[language=Python]
import statistics, math

def lag1_autocorr(xs):
    n = len(xs)
    m = statistics.mean(xs)
    num = sum((xs[i] - m) * (xs[i-1] - m) for i in range(1, n))
    den = sum((x - m) ** 2 for x in xs)
    return num / den

returns = [0.005, -0.002, 0.003, 0.001, ...]   # 対数リターン
rho = lag1_autocorr(returns)
se = 1 / math.sqrt(len(returns))
print(rho, se)   # |rho| > 2*se なら有意水準 5% で「記憶あり」
\end{lstlisting}

\begin{quote}
\textbf{【演習】}\ 500 日のリターンでラグ 1 自己相関 0.12 が得られた。有意か(標準誤差 \(1/\sqrt{500} \approx 0.045\))。また、有意でも「儲かる」までの距離について考察せよ。
\end{quote}

\textbf{解答の指針:}\ 0.12 は 0.09 を超えるので 5\% 水準で有意。しかし相関 0.12 は説明力 \(R^2 \approx 1.4\%\) に過ぎず、コストを引くと儲けに繋がるかは別問題。\textbf{統計的有意 ≠ 実務的有意}。

\section{AR モデル — 記憶をもつ系列の最小モデル}

\textbf{自己回帰モデル}( AR )は「昨日の自分で今日を説明する」モデルです。AR(1):
\[ X_t = c + \phi X_{t-1} + \varepsilon_t \]
\(\phi\) は持続性の強さ:\(|\phi| < 1\) で定常(衝撃は \(\phi^k\) に減衰して消える)、\(\phi = 1\) がランダムウォーク、\(|\phi| > 1\) は発散。

\begin{quote}
\textbf{【例】}\ \(\phi = 0.7\) の AR(1)。今日 +2\% のショックがあったら、明日は \(0.7 \times 2 = 1.4\%\) の押し上げ期待、明後日は 0.98\%、…と\textbf{衝撃が指数的に減衰して残る}。この「衝撃の残り方」が AR の本質です。
\end{quote}

次数を上げると AR(p): \(p\) 期前までを説明変数にします。次数の選定には ACF(と偏自己相関)の形を使います。\textbf{日次株価リターンは AR がほぼ効かない}(自己相関が僅か)ことで有名で、効きやすいのは\textbf{ボラティリティ}(大きい日は次の日も大きい)—— volatility clustering のモデル化が GARCH の発展に繋がります。

\begin{quote}
\textbf{【演習】}\ AR(1) で \(\phi = 1.05\) の系列の長期的挙動を述べよ。また現実のデータでこのような系列が非定常として排除される理由を述べよ。
\end{quote}

\textbf{解答の指針:}\ 衝撃が \(1.05^k\) で増幅され発散する。金額データ(人口・名目 GDP)がこうなりがちで、そのままモデル化すると見かけの強い相関が出るため、対数をとる・差分をとるなどの変換が必須です。

\section{見かけの回帰 — 非定常同士の罠}

トレンドを持つ 2 つの無関係な非定常系列(たとえば日本の人口とコーヒー消費量)を回帰すると、\textbf{高い相関と「有意」な回帰}が出ます。これが\textbf{見かけの回帰}( spurious regression )で、両系列が共通のトレンドを持つだけだからです。

防御法:\textbf{(1)} まず系列を定常にする(差分をとる)、\textbf{(2)} 層別に散布図を見る、\textbf{(3)} 因果のメカニズムを言葉で説明できるか考える(『統計の誤用事例集』第 1 話の再登場)。

\begin{quote}
\textbf{【演習】}\ 「水泳大会の参加者数」と「アイスの売上」の月次データに強い正の相関があった。時系列分析の観点からの防御手順を 3 ステップで述べよ。
\end{quote}

\textbf{解答の指針:}\ (1) 季節性があるなら差分や季節調整で定常化。(2) 年ごと・月ごとに層分けして相関の有無を再確認。(3) 共通要因(夏の暑さ)の説明メカニズムを検討する。

\section{予測の誠実さ — バックテストの作法}

時系列モデルを評価するときの作法は『機械学習入門』の過学習対策と同じ骨格ですが、\textbf{時間の順序を壊さない}点が特殊です:

\begin{enumerate}
\item \textbf{ウォークフォワード検証}: 過去で学習 → その直後で予測 → 窓を進める。未来を学習に使うことは\textbf{絶対に}しない(look-ahead bias)
\item \textbf{アウトオブサンプル}: 最後の期間を「触らない検証データ」として封印する
\item \textbf{ベースライン比較}: 「明日 = 今日」というナイーブ予測(ランダムウォーク予測)に勝てなければモデルの意味がない
\end{enumerate}

\begin{quote}
\textbf{【演習】}\ 株価データをランダムに train/test に分割して回帰モデルを評価した。この手順の致命的な問題を述べよ。
\end{quote}

\textbf{解答の指針:}\ 時系列では隣接点がほぼ同じ値のため、ランダム分割は「テストデータが実質学習データに隣接する」状態になり、性能を過大評価する(リーケージ)。分割は時間順に行う。

\section{おわりに — 時間構造の目}

この教科書の骨格: 定常性(§2)がすべての前提、自己相関(§3)が記憶の測定器、ランダムウォーク検定(§4)が「予測不能性」の検証、AR モデル(§5)が記憶の最小モデル、見かけの回帰(§6)が最大の罠、そして誠実な評価作法(§7)。

『確率過程入門』の理論(ランダムウォーク・AR は同じモデルの理論面)、『投資の数学』の期待値と資金管理(予測がどんなに良くてもコストとリスク管理は別)と繋げると、価格データをめぐる知識が一つの体系にまとまります。

\begin{quote}
\textbf{【総合演習】}
\begin{enumerate}
\item 非定常な系列を定常化する 2 つの操作を挙げよ。
\item AR(1) で \(\phi = 0.5\) のとき、今日の +4\% のショックの明日・明後日の押し上げ期待を求めよ。
\item ナイーブ予測(明日 = 今日)に勝てないモデルの意味を 1 文で述べよ。
\end{enumerate}
\end{quote}

\textbf{解答の指針:}\ 1. 差分をとる・対数をとる(または季節調整)。2. 明日 \(0.5 \times 4 = 2\%\), 明後日 1\%。3. そのデータには「予測に使える時間構造」がない(ランダムウォークであることと整合的)。

\end{document}
